% Options for packages loaded elsewhere
\PassOptionsToPackage{unicode}{hyperref}
\PassOptionsToPackage{hyphens}{url}
%
\documentclass[
  11pt,
]{article}
\usepackage{amsmath,amssymb}
\usepackage{iftex}
\ifPDFTeX
  \usepackage[T1]{fontenc}
  \usepackage[utf8]{inputenc}
  \usepackage{textcomp} % provide euro and other symbols
\else % if luatex or xetex
  \usepackage{unicode-math} % this also loads fontspec
  \defaultfontfeatures{Scale=MatchLowercase}
  \defaultfontfeatures[\rmfamily]{Ligatures=TeX,Scale=1}
\fi
\usepackage{lmodern}
\ifPDFTeX\else
  % xetex/luatex font selection
\fi
% Use upquote if available, for straight quotes in verbatim environments
\IfFileExists{upquote.sty}{\usepackage{upquote}}{}
\IfFileExists{microtype.sty}{% use microtype if available
  \usepackage[]{microtype}
  \UseMicrotypeSet[protrusion]{basicmath} % disable protrusion for tt fonts
}{}
\makeatletter
\@ifundefined{KOMAClassName}{% if non-KOMA class
  \IfFileExists{parskip.sty}{%
    \usepackage{parskip}
  }{% else
    \setlength{\parindent}{0pt}
    \setlength{\parskip}{6pt plus 2pt minus 1pt}}
}{% if KOMA class
  \KOMAoptions{parskip=half}}
\makeatother
\usepackage{xcolor}
\usepackage[margin=1in]{geometry}
\usepackage{color}
\usepackage{fancyvrb}
\newcommand{\VerbBar}{|}
\newcommand{\VERB}{\Verb[commandchars=\\\{\}]}
\DefineVerbatimEnvironment{Highlighting}{Verbatim}{commandchars=\\\{\}}
% Add ',fontsize=\small' for more characters per line
\newenvironment{Shaded}{}{}
\newcommand{\AlertTok}[1]{\textcolor[rgb]{1.00,0.00,0.00}{\textbf{#1}}}
\newcommand{\AnnotationTok}[1]{\textcolor[rgb]{0.38,0.63,0.69}{\textbf{\textit{#1}}}}
\newcommand{\AttributeTok}[1]{\textcolor[rgb]{0.49,0.56,0.16}{#1}}
\newcommand{\BaseNTok}[1]{\textcolor[rgb]{0.25,0.63,0.44}{#1}}
\newcommand{\BuiltInTok}[1]{\textcolor[rgb]{0.00,0.50,0.00}{#1}}
\newcommand{\CharTok}[1]{\textcolor[rgb]{0.25,0.44,0.63}{#1}}
\newcommand{\CommentTok}[1]{\textcolor[rgb]{0.38,0.63,0.69}{\textit{#1}}}
\newcommand{\CommentVarTok}[1]{\textcolor[rgb]{0.38,0.63,0.69}{\textbf{\textit{#1}}}}
\newcommand{\ConstantTok}[1]{\textcolor[rgb]{0.53,0.00,0.00}{#1}}
\newcommand{\ControlFlowTok}[1]{\textcolor[rgb]{0.00,0.44,0.13}{\textbf{#1}}}
\newcommand{\DataTypeTok}[1]{\textcolor[rgb]{0.56,0.13,0.00}{#1}}
\newcommand{\DecValTok}[1]{\textcolor[rgb]{0.25,0.63,0.44}{#1}}
\newcommand{\DocumentationTok}[1]{\textcolor[rgb]{0.73,0.13,0.13}{\textit{#1}}}
\newcommand{\ErrorTok}[1]{\textcolor[rgb]{1.00,0.00,0.00}{\textbf{#1}}}
\newcommand{\ExtensionTok}[1]{#1}
\newcommand{\FloatTok}[1]{\textcolor[rgb]{0.25,0.63,0.44}{#1}}
\newcommand{\FunctionTok}[1]{\textcolor[rgb]{0.02,0.16,0.49}{#1}}
\newcommand{\ImportTok}[1]{\textcolor[rgb]{0.00,0.50,0.00}{\textbf{#1}}}
\newcommand{\InformationTok}[1]{\textcolor[rgb]{0.38,0.63,0.69}{\textbf{\textit{#1}}}}
\newcommand{\KeywordTok}[1]{\textcolor[rgb]{0.00,0.44,0.13}{\textbf{#1}}}
\newcommand{\NormalTok}[1]{#1}
\newcommand{\OperatorTok}[1]{\textcolor[rgb]{0.40,0.40,0.40}{#1}}
\newcommand{\OtherTok}[1]{\textcolor[rgb]{0.00,0.44,0.13}{#1}}
\newcommand{\PreprocessorTok}[1]{\textcolor[rgb]{0.74,0.48,0.00}{#1}}
\newcommand{\RegionMarkerTok}[1]{#1}
\newcommand{\SpecialCharTok}[1]{\textcolor[rgb]{0.25,0.44,0.63}{#1}}
\newcommand{\SpecialStringTok}[1]{\textcolor[rgb]{0.73,0.40,0.53}{#1}}
\newcommand{\StringTok}[1]{\textcolor[rgb]{0.25,0.44,0.63}{#1}}
\newcommand{\VariableTok}[1]{\textcolor[rgb]{0.10,0.09,0.49}{#1}}
\newcommand{\VerbatimStringTok}[1]{\textcolor[rgb]{0.25,0.44,0.63}{#1}}
\newcommand{\WarningTok}[1]{\textcolor[rgb]{0.38,0.63,0.69}{\textbf{\textit{#1}}}}
\setlength{\emergencystretch}{3em} % prevent overfull lines
\providecommand{\tightlist}{%
  \setlength{\itemsep}{0pt}\setlength{\parskip}{0pt}}
\setcounter{secnumdepth}{-\maxdimen} % remove section numbering
\ifLuaTeX
  \usepackage{selnolig}  % disable illegal ligatures
\fi
\usepackage{bookmark}
\IfFileExists{xurl.sty}{\usepackage{xurl}}{} % add URL line breaks if available
\urlstyle{same}
\hypersetup{
  pdftitle={Full-family continuity from finite harmonic skeletons: a partial result for P03.3},
  pdfauthor={Fresh independent draft},
  hidelinks,
  pdfcreator={LaTeX via pandoc}}

\title{Full-family continuity from finite harmonic skeletons: a partial
result for P03.3}
\author{Fresh independent draft}
\date{2026-10-04}

\begin{document}
\maketitle

\section{Result and exact scope}\label{result-and-exact-scope}

The sharp classification requested in P03.3 is \textbf{not proved here}.
In particular, this draft does not assert continuity on every normal or
geometrically unibranch analytic domain. It proves a finite-graph
obstacle theorem, a compact-skeleton continuity criterion for an entire
competitor class, and a proper-fibre reduction. All three proofs are
given below. The criterion resolves a specific uniformity/cofinality
mechanism: finite-skeleton envelopes dominate \textbf{every} competitor
by comparison, and uniform approximation of the obstacle then controls
the full envelope. No modulus of continuity for the competitors is
assumed.

There is an important application gap. This draft does not verify all
the criterion's hypotheses for the exact locally
Fubini--Study/semipositive-real-PL union on any new positive-dimensional
analytic class. Rational affinoids in the analytic projective line and
their products with proper projective factors are proposed applications,
not established geometric continuity theorems. The supplied unpublished
claims, audit verdicts, and hashes are not used as proof inputs. No
theorem about classical psh or CLD-v2 membership of the envelope is
claimed.

The graph and compact-skeleton statements are mathematical proofs,
subject to independent checking. The small rational computation at the
end is a check of the graph construction, not a proof of an analytic
identification or a classification.

\section{1. A finite graph obstacle
problem}\label{a-finite-graph-obstacle-problem}

Let \(\Gamma\) be a nonempty finite connected metric graph. Subdivide it
so that each edge joins distinct vertices; parallel edges are allowed.
Write \(V\) for its vertices, \(E\) for its edges, and \(\ell_e>0\) for
the length of an edge. Fix a possibly empty set \(B\subset V\) of
designated boundary vertices. The vertex set may contain degree-two
subdivision vertices. Let \(g\) be a continuous function affine on each
edge, with vertex values \(g_v\).

For \(u\in\mathbb R^V\) define the outgoing-slope sum

\[
L_vu=\sum_{e=vw}\frac{u_w-u_v}{\ell_e}.
\]

Each parallel edge is included separately. Let \(\widetilde u\) be the
affine interpolation of \(u\) on \(\Gamma\). Define the feasible set

\[
\mathcal F_g=\{u\in\mathbb R^V:u_v\leq g_v\ (v\in V),\quad L_vu\geq0\ (v\notin B)\}.
\]

Thus an affine-on-edges function is graph-subharmonic when its outgoing
slopes have nonnegative sum at every nonboundary vertex. No inequality
is imposed at a designated boundary vertex.

\textbf{Proposition 1 (greatest graph minorant and complementarity).}
The set \(\mathcal F_g\) has a greatest element \(p\) in coordinatewise
order. Its interpolation satisfies \(\widetilde p\leq g\). At every
\(v\in B\),

\[p_v=g_v.\]

At every \(v\notin B\),

\[p_v<g_v\quad\Longrightarrow\quad L_vp=0.\]

Consequently \(\widetilde p\) is harmonic on the open set
\(\{\widetilde p<g\}\): it is affine on open edge segments and has zero
total outgoing slope at every vertex in that set. This open set misses
\(B\).

\textbf{Proof.} Put \(m=\min_{v\in V}g_v\). The constant vector \(m\) is
feasible. If \(u,w\in\mathcal F_g\), their coordinatewise maximum \(z\)
is feasible: at a vertex \(v\) choose \(u\) with \(z_v=u_v\)
(interchange \(u,w\) if necessary), and use \(z_t\geq u_t\) at every
neighboring vertex \(t\) to obtain \(L_vz\geq L_vu\geq0\). This also
proves that replacing any feasible vector by its maximum with \(m\)
preserves feasibility.

The nonempty set

\[\mathcal F_g\cap\prod_{v\in V}[m,g_v]\]

is compact. Choose a maximizer \(p\) of \(\sum_vu_v\) there. For any
feasible \(u\), the vector \(\max(p,u)\) belongs to this compact set. If
\(u_v>p_v\) at some coordinate, its sum is strictly larger. Therefore
\(u\leq p\); this proves the greatest-element assertion, without
assuming a priori that a selected family of feasible vectors is cofinal.

For \(v\in B\), increasing \(p_v\) up to \(g_v\) preserves every
constraint: there is no constraint at \(v\), and all neighboring
outgoing-slope sums increase. Hence \(p_v=g_v\). If \(v\notin B\),
\(p_v<g_v\), and \(L_vp>0\), increase just \(p_v\) by a sufficiently
small positive \(\delta\). Its own outgoing-slope sum decreases by

\[\delta\sum_{e\ni v}\ell_e^{-1},\]

and stays nonnegative; all neighboring sums increase. The obstacle
constraint stays valid. This contradicts maximality. If the graph has
only one vertex and no edges, its slope sum is already zero and the
conclusion follows directly. Finally, affine interpolation preserves the
vertex obstacle inequalities throughout every edge. The stated
harmonicity follows from complementarity. \(\square\)

\textbf{Proposition 2 (empty graph boundary).} If \(B=\varnothing\), the
greatest feasible vector is the constant \(m=\min_Vg_v\).

\textbf{Proof.} For every vector \(u\), \(\sum_vL_vu=0\), because the
two contributions from an edge cancel. Thus a feasible vector has
\(L_vu=0\) at every vertex. Choose a vertex at which \(u\) is maximal.
Every summand in its outgoing-slope sum is nonpositive, so every
neighboring value equals the maximal value. Connectivity propagates
equality to all vertices. Hence all feasible vectors are constant, and
the largest admissible constant is \(m\). \(\square\)

\subsection{An exact finite algorithm}\label{an-exact-finite-algorithm}

When \(B\ne\varnothing\), enumerate all subsets \(A\subset V\)
containing \(B\). Impose

\[u_v=g_v\quad(v\in A),\qquad L_vu=0\quad(v\notin A).\]

This is a nonsingular Dirichlet linear system: a solution of the
homogeneous system has zero values on \(A\) and is harmonic elsewhere;
propagation from a maximal and a minimal vertex, using connectivity,
makes it identically zero. Solve the system and retain precisely the
solutions satisfying all feasibility inequalities. Proposition 1 shows
that the true \(p\) appears: take \(A=\{v:p_v=g_v\}\), which contains
\(B\), and use complementarity. A retained vector that dominates all
retained vectors is therefore \(p\). With rational lengths and obstacle
values this procedure uses exact rational arithmetic. The number of
systems is at most \(2^{|V\setminus B|}\). This is a finite algorithm
for a supplied finite graph, not an algorithm for the original analytic
classification.

\section{2. A compact-skeleton criterion for a full competitor
class}\label{a-compact-skeleton-criterion-for-a-full-competitor-class}

The following hypotheses describe harmonic geometry and individual
competitors, not continuity of the desired envelope.

Let \(X\) be a nonempty compact connected Hausdorff space, and let
\(D\subset X\) be a finite, possibly empty, designated boundary set. A
harmonic function on an open subset of \(X\setminus D\) will be a member
of a specified collection \(\mathcal H\) of continuous functions; this
collection must be invariant under adding constants and restricting to
open subsets. Let \(\mathcal C\subset C(X,\mathbb R)\) be an
\textbf{entire} given competitor class.

A finite harmonic skeleton presentation for
\((X,D,\mathcal H,\mathcal C)\) consists of the following data and
hypotheses.

\begin{enumerate}
\def\labelenumi{\arabic{enumi}.}
\item
  There is a collection of continuous retractions \(r:X\to\Gamma\) onto
  finite connected metric graphs embedded in \(X\). Every \(\Gamma\)
  contains \(D\) as designated graph boundary vertices \(B=D\), and
  \(r\) fixes \(\Gamma\) pointwise. Subdivision at additional vertices
  is allowed. Graph-affine means affine on every edge of a chosen finite
  subdivision.
\item
  Every \(f\in C(X,\mathbb R)\) is a uniform limit of functions
  \(g=\gamma\circ r\), where \(\gamma\) is graph-affine on one of these
  finite subdivided skeletons. These approximating obstacles need not be
  competitors or subharmonic.
\item
  For every such skeleton and every graph-affine \(a\) with
  \(L_va\geq0\) for \(v\notin D\), its lift \(a\circ r\) belongs to
  \(\mathcal C\). All constant functions belong to \(\mathcal C\), and
  \(h+c\in\mathcal C\) whenever \(h\in\mathcal C\) and
  \(c\in\mathbb R\).
\item
  For every graph-affine \(a\), the lift \(a\circ r\) is harmonic on the
  inverse image of every graph-open set avoiding \(D\) on which \(a\)
  has zero outgoing-slope sum at each vertex. On an open edge segment,
  affine functions count as harmonic. In particular, constant lifts on
  the components discarded by the retraction must have this property.
  This is a substantive local harmonic-lifting hypothesis, including at
  branch attachments.
\item
  Every \(h\in\mathcal C\) has harmonic comparison on relatively compact
  open sets in \(X\setminus D\). Precisely: if \(W\) is open,
  \(\overline W\subset X\setminus D\), \(W\) has nonempty topological
  boundary \(\partial_XW\), \(q\) is harmonic on an open neighborhood of
  \(\overline W\), and \(h\leq q\) on \(\partial_XW\), then \(h\leq q\)
  on \(W\).
\end{enumerate}

No assertion that these hypotheses are automatic in the original
analytic setting is made. In particular, hypothesis 3 is a global
realization assertion and hypothesis 5 concerns \textbf{every} member of
the exact competitor union.

For a finite continuous obstacle \(f\) on \(X\), set

\[P_{\mathcal C}f(x)=\sup\{h(x):h\in\mathcal C,\ h\leq f\text{ on }X\}.\]

\textbf{Theorem 3 (full-family continuity).} Under hypotheses 1--5,
\(P_{\mathcal C}f\) is finite and continuous for every
\(f\in C(X,\mathbb R)\). It equals its upper-semicontinuous
regularization. For a skeletal obstacle \(g=\gamma\circ r\), it is
exactly \(\widetilde p\circ r\), where \(p\) is the greatest vector of
Proposition 1 for \(\gamma\). For arbitrary \(f\), if

\[\|g_j-f\|_\infty\leq\varepsilon_j\longrightarrow0,\]

then

\[\|P_{\mathcal C}f-\widetilde p_j\circ r_j\|_\infty\leq\varepsilon_j.\]

This is an assertion about the full class \(\mathcal C\), not only its
skeletal members.

\textbf{Proof.} Nonemptiness and finiteness follow from

\[\min_X f\leq P_{\mathcal C}f\leq f.\]

First take a skeletal obstacle \(g=\gamma\circ r\). Put
\(q=\widetilde p\circ r\). Hypothesis 3 and Proposition 1 give
\(q\in\mathcal C\) and \(q\leq g\), so \(q\leq P_{\mathcal C}g\).

Let \(h\) be \textbf{any} member of \(\mathcal C\) with \(h\leq g\).
Suppose \(h(x)>q(x)\) somewhere, and choose \(0<\eta<h(x)-q(x)\). Set

\[W=\{z:h(z)>q(z)+\eta\}.\]

This open set is nonempty, and its closure is compact. Throughout
\(\overline W\), continuity gives \(h\geq q+\eta\), hence
\(g-q\geq\eta\). Thus \(\overline W\) is contained in the open set
\(\{q<g\}\), where \(q\) is harmonic by Proposition 1 and hypothesis 4.
This open set misses \(D\), because \(p_v=\gamma_v\) at boundary
vertices.

We must also check that \(W\) has nonempty boundary, including when
\(D=\varnothing\). The contact set \(\{q=g\}\) is nonempty. If
\(D\ne\varnothing\), this follows from boundary contact. If
\(D=\varnothing\), Proposition 2 gives a constant \(q=\min\gamma\),
attained by \(\gamma\) at some graph vertex. At any contact point,
\(h\leq g=q\), so the point is outside \(\overline W\). Therefore \(W\)
is a proper nonempty open subset of connected \(X\). If it had empty
boundary, it would be both open and closed, contrary to connectivity.
Moreover \(h=q+\eta\) on \(\partial_XW\), by continuity and the
definition of \(W\).

Harmonic comparison with \(q+\eta\) therefore yields \(h\leq q+\eta\)
throughout \(W\), contradicting its definition. Hence \(h\leq q\) on
\(X\). Since this argument applies to every competitor,
\(P_{\mathcal C}g=q\).

Now note that the full-envelope operator has uniform stability. If
\(\|f-g\|_\infty\leq\varepsilon\) and \(h\in\mathcal C\) satisfies
\(h\leq f\), then \(h-\varepsilon\in\mathcal C\) satisfies
\(h-\varepsilon\leq g\). Taking the supremum gives
\(P_{\mathcal C}f\leq P_{\mathcal C}g+\varepsilon\). Interchanging
\(f,g\) gives

\[\|P_{\mathcal C}f-P_{\mathcal C}g\|_\infty\leq\varepsilon.\]

Apply this to the obstacles from hypothesis 2. Their exact full
envelopes are the continuous functions \(\widetilde p_j\circ r_j\) just
proved above. They converge uniformly to \(P_{\mathcal C}f\). A uniform
limit of continuous functions is continuous on any topological space:
given a point and a positive tolerance, one fixed approximant and its
local continuity provide the required neighborhood. For a continuous
function \(u\), the neighborhood formula

\[u^*(x)=\inf_{N\ni x}\sup_Nu\]

equals \(u(x)\). Thus regularization leaves \(P_{\mathcal C}f\)
unchanged. \(\square\)

\subsection{What cofinality was actually
proved}\label{what-cofinality-was-actually-proved}

The majorization step is essential. The fact that graph minorants are
continuous, or that some sequence of them has a continuous limit, would
not bound non-skeletal competitors. In the proof, every competitor
\(h\leq g_j\) is bounded by \(q_j=P_{\mathcal C}g_j\). Moreover
\(q_j-\varepsilon_j\) is a \textbf{global} competitor below \(f\), since

\[q_j-\varepsilon_j\leq g_j-\varepsilon_j\leq f.\]

For any arbitrary global competitor \(h\leq f\), the translated function
\(h-\varepsilon_j\) lies below \(g_j\), so

\[q_j-\varepsilon_j\leq P_{\mathcal C}f\leq q_j+\varepsilon_j.\]

These two inequalities close the full-family gap; no assumption that the
selected sequence was exhaustive is used. The construction does not
require maxima of arbitrary members of \(\mathcal C\) to remain in
\(\mathcal C\).

If \(D=\varnothing\), every skeletal envelope is the constant minimum of
its obstacle by Proposition 2. Uniform stability then gives

\[P_{\mathcal C}f\equiv\min_Xf.\]

In that case every \(h\in\mathcal C\) is constant: apply this formula to
the obstacle \(f=h\), for which \(h\) is itself an admissible
competitor. This is a consequence of all the skeleton hypotheses, not an
independently proved Liouville theorem for general proper analytic
spaces.

\subsection{The criterion does not imply a common competitor
modulus}\label{the-criterion-does-not-imply-a-common-competitor-modulus}

Take \(X=[0,1]\), designate both endpoints as boundary, take ordinary
affine harmonic functions locally, and let \(\mathcal C\) be the class
of finite continuous convex functions. The hypotheses hold, using
interval subdivisions. Harmonic comparison is elementary: on each
component interval, a convex function below an affine function at its
endpoints remains below it throughout by the chord inequality. The
functions

\[h_n(t)=\max\{0,1-nt\},\qquad n\geq1,\]

all lie below the continuous obstacle \(1\). They have \(h_n(0)=1\) and
\(h_n(1/n)=0\), so they do not share a modulus of continuity at \(0\).
Nevertheless the envelope over \textbf{all} finite continuous convex
minorants of \(1\) is the continuous constant \(1\), since that constant
is itself a competitor. Theorem 3 therefore uses a different mechanism
from a common local modulus for the full family.

\section{3. How to prove the obstacle-density
hypothesis}\label{how-to-prove-the-obstacle-density-hypothesis}

For use in an application, hypothesis 2 can be split into separation and
a mesh statement. This section supplies the elementary approximation
argument instead of citing a density theorem without checking its
hypotheses.

\textbf{Lemma 4 (lattice approximation).} Let \(X\) be compact
Hausdorff. Let \(\mathcal A\subset C(X,\mathbb R)\) contain constants,
be closed under real linear combinations, finite maxima and finite
minima, and separate distinct points. Then every continuous function on
\(X\) is uniformly approximable by elements of \(\mathcal A\).

\textbf{Proof.} Fix \(f\in C(X,\mathbb R)\) and \(\epsilon>0\). For each
pair \(x,y\), separation and affine rescaling give
\(a_{x,y}\in\mathcal A\) with \(a_{x,y}(x)=f(x)\) and
\(a_{x,y}(y)=f(y)\); for \(x=y\) use a constant. For a fixed \(x\), the
open sets

\[\{z:a_{x,y}(z)>f(z)-\epsilon\},\qquad y\in X,\]

cover \(X\). Choose finitely many \(y\) and let \(b_x\) be the maximum
of their functions. Then \(b_x>f-\epsilon\) on \(X\) and
\(b_x(x)=f(x)\). The open sets \(\{z:b_x(z)<f(z)+\epsilon\}\), as \(x\)
varies, also cover \(X\). Choose finitely many \(x\) and let \(a\) be
the minimum of the corresponding \(b_x\). It satisfies
\(f-\epsilon<a<f+\epsilon\) on \(X\) and belongs to \(\mathcal A\).
\(\square\)

Suppose that continuous functions factoring through the finite skeleta
form such a lattice, because a common refinement absorbs each finite
collection of skeleta and their functions. Suppose also that any
continuous function on one finite graph can be uniformly approximated by
affine interpolation on allowable subdivisions. The latter holds when
allowed subdivision points are dense on every edge: choose a
sufficiently fine subdivision using uniform continuity on the finite
graph. Lemma 4 and interpolation then prove hypothesis 2. A common
refinement must be checked in the specific application; it is not
supplied by compactness alone. Neither separation nor interpolation
checks the semipositivity or global realization hypothesis 3.

\section{4. Proper fibres: reduction of the whole class, not
symmetrization}\label{proper-fibres-reduction-of-the-whole-class-not-symmetrization}

This is a separate exact transfer theorem. It makes clear both the
strength of the proper-fibre approach and the missing analytic
obligations.

Let \(p:Y\to X\) be a continuous, closed, open surjection of Hausdorff
spaces with nonempty compact fibres. Suppose it has a continuous section
\(s:X\to Y\). Let \(\mathcal C_X\subset C(X,\mathbb R)\) and
\(\mathcal C_Y\subset C(Y,\mathbb R)\) be competitor classes. Assume:

\begin{itemize}
\tightlist
\item
  Pullback along \(p\) sends \(\mathcal C_X\) into \(\mathcal C_Y\).
\item
  Pullback along \(s\) sends \(\mathcal C_Y\) into \(\mathcal C_X\).
\item
  Every member of \(\mathcal C_Y\) is constant on each fibre of \(p\).
\end{itemize}

The last condition could be proved through stability of the exact class
under analytic fibre restriction and a fibre Liouville theorem. Neither
of those analytic assertions is assumed proved in this draft.

\textbf{Theorem 5 (arbitrary-obstacle fibre reduction).} For every
finite continuous \(f:Y\to\mathbb R\), the function

\[g(x)=\min_{y\in p^{-1}(x)}f(y)\]

is finite and continuous, and

\[P_{\mathcal C_Y}f=P_{\mathcal C_X}g\circ p.\]

This identity includes the empty-family case with supremum \(-\infty\).
In the finite nonempty case, continuity of the base envelope implies
continuity of the full envelope upstairs. All competitors are used; no
invariant subfamily is substituted.

\textbf{Proof.} Compactness of a fibre and continuity of \(f\) give its
finite attained minimum. For any \(t\in\mathbb R\),

\[\{x:g(x)<t\}=p(\{y:f(y)<t\})\]

is open, since \(p\) is open. Also

\[\{x:g(x)\leq t\}=p(\{y:f(y)\leq t\})\]

is closed, since \(p\) is closed; attainment of the minimum justifies
this equality. Thus \(g\) is both upper and lower semicontinuous, hence
continuous.

For any \(h\in\mathcal C_Y\), put \(a=h\circ s\in\mathcal C_X\). Fibre
constancy implies \(h=a\circ p\). Moreover \(h\leq f\) if and only if
\(a\leq g\). Conversely, if \(a\in\mathcal C_X\) and \(a\leq g\), its
pullback belongs to \(\mathcal C_Y\) and lies below \(f\). Hence the two
admissible families are in exact bijection under pullback. Taking
pointwise suprema proves the formula, including simultaneous emptiness.
\(\square\)

For compact connected \(X\) satisfying Theorem 3, and for the full pair
of classes satisfying these fibre hypotheses, Theorems 3 and 5 give
continuity on \(Y\) for \textbf{all} finite continuous obstacles. This
is a complete proof of that conditional transfer statement. Openness in
Theorem 5 is not removable from the proof of upper semicontinuity of
\(g\) merely because the fibres are compact. Closedness alone supplies
only lower semicontinuity.

\section{5. Exact obligations for an analytic
application}\label{exact-obligations-for-an-analytic-application}

The following are genuine unchecked analytic gates, rather than claimed
facts extracted from the supplied dossier.

\subsection{Rational curve domains}\label{rational-curve-domains}

A prospective base class is a connected compact strict rational affinoid
in \((\mathbb P^1_K)^{\mathrm{an}}\), over a complete algebraically
closed nontrivially valued field \(K\). The proposed boundary set is its
relative affinoid boundary. Finite trees obtained by adding finitely
many disc points to a boundary skeleton are the proposed graphs.

To turn Theorem 3 into a theorem for the original competitor union on
this class, one must establish all of the following under the literal
analytic definitions:

\begin{enumerate}
\def\labelenumi{\arabic{enumi}.}
\item
  The indicated retractions are continuous, their boundary vertices are
  exactly appropriate relative-boundary points, and the
  density/refinement assertions of Section 3 hold. Type-1 endpoints,
  nondiscrete radii, and type-4 points cannot be omitted.
\item
  Every finite graph-affine function with nonnegative interior outgoing
  slopes lifts to a \textbf{global semipositive real-PL function} on the
  entire affinoid domain, with the required permitted vertices and
  coefficients. It is not enough to produce a continuous
  curvewise-subharmonic function, a local expression, or an approximant
  that misses the obstacle. This is the globalization/realization gate.
\item
  Every locally Fubini--Study competitor and every semipositive real-PL
  competitor satisfies the harmonic comparison used in hypothesis 5. The
  signs and weights of the slope operator must agree with the precise
  analytic convention. Interior branch attachments have to be included,
  while relative-boundary vertices must have exactly the asserted
  freedom. A comparison result for only one branch of the union would be
  insufficient.
\item
  Lifted graph-harmonic functions are harmonic on the required open
  inverse images, including discarded disc components and branch
  attachment points. Continuity of the retraction alone does not prove
  this.
\end{enumerate}

If these gates are supplied, the full-family upper bound and the global
minorants are then furnished by Theorem 3, rather than assumed. This
draft has not independently completed gates 1--4 or obtained a primary
text verifying their exact applicability. Consequently it does not
assert even this rational-affinoid application as established P03.3
progress on a new analytic class.

\subsection{Scalar extensions}\label{scalar-extensions}

An analytic application must separately check the gates after scalar
extension and verify how the skeletal data and \textbf{the exact
competitor class} pull back. A useful conditional statement is the
following: suppose \(\mathcal G\subset C(X,\mathbb R)\) is a uniformly
dense collection of skeletal obstacles. For a continuous scalar
projection \(\pi:X'\to X\), assume that every \(g\in\mathcal G\) has a
finite exact full envelope \(q=P_{\mathcal C}g\) and the exact full
envelope upstairs satisfies

\[P_{\mathcal C'}(g\circ\pi)=q\circ\pi,\]

with both classes invariant under adding constants. Then for every
finite continuous \(f\) on compact \(X\),

\[P_{\mathcal C'}(f\circ\pi)=P_{\mathcal C}f\circ\pi.\]

Indeed, choose \(g_j\in\mathcal G\) with
\(\|f-g_j\|_\infty\leq\varepsilon_j\to0\), and put
\(q_j=P_{\mathcal C}g_j\). The same approximation bound holds after
composition with \(\pi\), and uniform stability bounds both sides within
\(\varepsilon_j\) of \(q_j\circ\pi\). The finite skeletal envelope
identities ensure nonempty admissible families; translation supplies a
minorant for each nearby obstacle as well. Let \(j\) tend to infinity.
This argument does not require an open scalar projection. Its skeletal
pullback identity is nevertheless an additional full-family assertion
that must be proved; it cannot be inferred from pullback stability
alone. Nothing here proves identities between upper regularizations
before continuity has been established.

\subsection{Higher dimensional
products}\label{higher-dimensional-products}

A proposed higher dimensional range is a rational curve affinoid
multiplied by a finite product of analytic projective lines. To apply
Theorem 5 to the projection, one must prove:

\begin{itemize}
\tightlist
\item
  Restriction and pullback stability of the literal full
  locally-Fubini--Study/semipositive-real-PL union, including
  restrictions to fibres over completed residue fields.
\item
  A Liouville theorem for both branches of that union on the proper
  fibres, after every necessary field extension, with connectedness and
  branch compatibility checked.
\item
  Closedness, compact fibres, and \textbf{openness} of the analytic
  projection, together with an analytic section providing the
  class-preserving restriction in Theorem 5.
\item
  The curve-base gates above for the continuous fibre-minimum obstacle,
  not merely for invariant or pulled-back original obstacles.
\end{itemize}

These checks have not been completed. Properness by itself does not
prove the fibre Liouville statement for an arbitrary locally defined
class. Nor does a global section by itself prove constancy on a fibre or
openness at its other points.

\subsection{What remains for the literal original
problem}\label{what-remains-for-the-literal-original-problem}

Even completion of the proposed curve and product applications would
give only a sufficient geometric sector. The original target
additionally requires sharp hypotheses and necessity across the residual
higher dimensional normal/geometrically unibranch range, with arbitrary
finite continuous obstacles, scalar extension, branches, relative
boundary, and global competitor issues. No necessity theorem or
counterexample in that residual range is proved here. Positivity of an
envelope would not replace the continuity argument, and continuity
obtained here would not independently imply the requested classical or
CLD-v2 positivity statements.

\section{6. Reproducibility and
evidence}\label{reproducibility-and-evidence}

The authoritative fresh input is preserved byte-for-byte as
\texttt{input.json} and \texttt{input-original.json};
\texttt{task-original.json} and \texttt{assignment-original.txt} were
written only because they were absent. The mathematical input's
\texttt{{[}OCR\ control\ U+XXXX{]}} markers were not given a
mathematical interpretation.

The exact-arithmetic program is \texttt{finite\_graph\_envelope.py}.

Its input is \texttt{finite-graph-input.json}. It writes
\texttt{finite-graph-output.json}. Run

\begin{Shaded}
\begin{Highlighting}[]
\NormalTok{python finite\_graph\_envelope.py \textbackslash{}}
\NormalTok{  {-}{-}input finite{-}graph{-}input.json \textbackslash{}}
\NormalTok{  {-}{-}output finite{-}graph{-}output.json}
\end{Highlighting}
\end{Shaded}

It enumerates Dirichlet contact sets and solves using rational Gaussian
elimination. It checks feasibility, domination of all enumerated
feasible contact solutions, contact/harmonic complementarity, and
translation stability in the recorded examples. These are exhaustive
contact-set computations for those supplied finite graphs only. They do
not certify any analytic comparison or model-metric identification.

The exact computed greatest vertex vectors, in the input's vertex order,
are:

\begin{itemize}
\tightlist
\item
  A three-vertex path with boundary at both ends and obstacle
  \((0,10,0)\): \((0,0,0)\).
\item
  A four-vertex path with boundary at both ends and obstacle
  \((4,3,0,4)\): \((4,2,0,4)\).
\item
  A rooted three-vertex path with obstacle \((2,0,2)\): \((2,0,0)\).
\item
  A rooted four-vertex star with obstacle \((3,0,2,4)\): \((3,0,2,3)\).
\item
  A four-vertex cycle with one boundary vertex and obstacle
  \((3,10,0,10)\): \((3,3/2,0,3/2)\).
\item
  A connected four-vertex path with empty designated boundary and
  obstacle \((5,-2,7,1)\): \((-2,-2,-2,-2)\).
\item
  A single vertex with empty designated boundary and obstacle \(3/7\):
  \(3/7\).
\item
  A three-vertex graph with two parallel edges of lengths \(1,2\)
  between the first two vertices, an edge of length \(1\) to the third,
  both endpoints designated boundary, and obstacle \((4,10,0)\):
  \((4,12/5,0)\).
\end{itemize}

Public-source provenance is in \texttt{source-provenance.json}. Direct
HTTP retrieval of arXiv and an author's public page failed through the
unchanged existing proxy. Public GitHub discovery did not supply a
primary mathematical text for the analytic gates. No retrieved primary
theorem is used. Literal source extracts therefore have no content;
\texttt{source-extracts/README.md} records that absence. Search snippets
and the user dossier are not substituted for primary proofs.

The standalone TeX paper is \texttt{P03.3-paper.tex}. It contains this
entire proof and scope statement; it is generated from this Markdown
with \texttt{render-paper.sh}. Actual twice-run TeX commands and
diagnostics are recorded by \texttt{build-paper.py} in
\texttt{build-record.json} and separate log files. The rerun sequence is

\begin{Shaded}
\begin{Highlighting}[]
\NormalTok{bash render{-}paper.sh}
\NormalTok{python build{-}paper.py}
\NormalTok{python verify{-}artifacts.py}
\end{Highlighting}
\end{Shaded}

\texttt{environment.json}, \texttt{commands.json},
\texttt{verification.json}, and \texttt{status.json} record execution
details and limitations. No installation, environment-setting change,
other model invocation, or supervised large computation was used. There
is no resource-heavy calculation requiring a supervised compute
facility. Successful compilation, hashes, and finite tests establish
only the stated execution facts.

\end{document}
